% 第 12 章 Critical Area Analysis Flow
\bichapter{关键面积分析流程}{Critical Area Analysis Flow}
\label{chap:caa}

本章介绍 IC Validator 关键面积分析（critical area analysis）流程。

IC Validator 关键面积分析（CAA）流程用于识别版图中对制造环境中随机污染物所致失效敏感的区域。
关键面积定义为：给定尺寸缺陷的中心若落在该区域内，可能导致电路短路或开路失效。
CAA 流程当前支持两类缺陷：OPEN 与 SHORT。

\figplaceholder{Figure 90: Critical Area by Defect Types}{按缺陷类型划分的关键面积}{fig:caa-defect-types}

关键面积分析流程可用下列示意图说明：

\figplaceholder{Figure 91: Critical Area Analysis Flow}{关键面积分析流程}{fig:caa-flow}

\begin{enumerate}
  \item \textbf{晶圆厂缺陷数据：} CAA 流程的起点是晶圆厂提供的制造缺陷数据及缺陷尺寸计算公式。\\
  缺陷数据通常包含由三个参数定义的缺陷尺寸范围——最小与最大缺陷尺寸 $D_{\mathrm{min}}$、$D_{\mathrm{max}}$、缺陷尺寸采样数，以及在 $(D_{\mathrm{min}}, D_{\mathrm{max}})$ 范围内计算缺陷尺寸的系数。

  \item \textbf{查找版图关键面积：} 使用 \cmd{critical\_areas()} 函数确定版图关键面积。更多信息见 \textit{IC Validator Reference Manual} 中的 \cmd{critical\_areas()} 函数。

  \item \textbf{分析关键面积/生成报告：} 识别出版图关键面积并以多边形层表示后，可在 runset 中用这些数据计算加权关键面积（WCA）与平均质量（average quality）。\\
  WCA 计算为
  \[
    \mathrm{WCA} = K \sum_{i} \frac{\mathrm{CA}_{i}}{D_{i}^{\mathrm{fall\_power}}}
  \]
  其中 K-factor 由用户在 runset 中提供。\\
  平均质量通常按如下计算：
  \[
    \mathrm{Average\ Quality} = 1 - \frac{\mathrm{WCA}}{\mathrm{AREA}}
  \]
  其中 WCA 为给定版图层的加权关键面积，AREA 为该层总面积。\\
  计算结果可用于生成报告与热图。热图应使用 VUE 工具生成。\\
  最简单的关键面积报告可由 \cmd{dfm\_feature()} 函数生成，并输出到所分析层的 \cmd{LAYOUT\_ERRORS} 文件。
\begin{lstlisting}[style=shell]
chip area:                         = 1543.3206
Critical_Area[0]:                      = 0.0000
Defect Size[0]:                        = 0.0110
Critical_Area[1]:                      = 0.0000
Defect Size[1]:                        = 0.0120
Critical_Area[2]:                      = 0.0000
Defect Size[2]:                        = 0.0190
Critical_Area[3]:                      = 0.0000
Defect Size[3]:                        = 0.0380
Critical_Area[4]:                      = 0.0000
Defect Size[4]:                        = 0.0740
Critical_Area[5]:                      = 0.0000
Defect Size[5]:                        = 0.1350
Critical_Area[6]:                      = 1.4728
Defect Size[6]:                        = 0.2250
Critical_Area[7]:                      = 4.9703
Defect Size[7]:                        = 0.3500
Critical_Area[8]:                      = 17.1042
Defect Size[8]:                        = 0.5170
Critical_Area[9]:                      = 90.4452
Defect Size[9]:                        = 0.7320
Critical_Area[10]:                     = 200.9118
Defect Size[10]:                       = 1.0000
WCA:                                   = 0.0005
DIFF_CA_OPEN AVERAGE QUALITY:          = 1.0000
\end{lstlisting}
  关键面积热图由 VUE 工具生成，如图~\ref{fig:caa-heatmap} 所示。
\end{enumerate}

\figplaceholder{Figure 92: Heat Maps in VUE}{VUE 中的热图}{fig:caa-heatmap}

% ============================================================
\section{IC Validator 关键面积分析器 Runset}
\subsection*{IC Validator Critical Area Analyzer Runset}

下列 PXL 结构给出在 CAA runset 中组织数据的一种可能方式：
\begin{lstlisting}[style=pxl]
ca_defect_range_s : newtype struct of {
   count : integer;
   min : double;
   max : double;
};
ca_defect_density_s : newtype struct of {
   na : integer;            // dummy
   fall_power : double;
   conversion_factor : double;
};
ca_config_s : newtype struct of {
   layer : polygon_layer;
   mode : critical_area_mode_e;
   modeling : integer     // dummy
   range : ca_defect_range_s;
   defect_density : ca_defect_density_s;
};
ca_layer_configuration_l : newtype list of ca_config_s;
\end{lstlisting}

典型 CAA runset 如下：
\begin{lstlisting}[style=pxl]
#include <icv.rh>
library("input_desing.gds", GDSII, "top");

// assign section
DIFF = assign({ { 3 } });
POLY = assign({ { 3 } });
METAL1 = assign({ { 3 } });

connect_layer = assign({ { 4 } });

// types to store defect sizes;
defect_sizes_list : newtype list of double;
defect_sizes_h : newtype hash of integer to defect_sizes_list;
defect_hash : defect_sizes_h = {};
layer_dr : list of double = {};
layer_tagname : string = "";
fall_power : double = 0.0;
ca_mode : critical_area_mode_e;

g_ca_layers : string_to_geometry_layer_h = {};

// the remote function to calculate the WCA and average quality and
// output report
get_area : function (void) returning void
{
     dfm_report_double("chip area: ", dfm_window_area());

     wca_tmp: double = 0;

     // calculate K-factor
     // put actual K-calculation code here
     K = 6.0e-07;

     // for each input layer the critical_areas() command creates
     // list of critical area layers - one layer per defect size

    for (idx = 0 to layer_dr.size() - 1) {
      this_area = dfm_sum_area(""+idx);
      dfm_report_double("Critical_Area["+idx+"]: ", this_area);
      // output critical area value

      dfm_report_double("Defect Size["+idx+"]: ", layer_dr[idx]);
      // output defect size

         // calculate sum of WCAs for all defect sizes for given input layer
         wca_tmp= wca_tmp + (this_area / (layer_dr[idx]**fall_power));
     }

     // final WCA calculation
     wca = K * wca_tmp;

     dfm_report_double("WCA: " , wca);        // report WCA

    average_quality : double = 1 - (wca / dfm_window_area());
    // calculate average quality
    dfm_report_double(layer_tagname + "_" + ca_mode + " AVERAGE QUALITY: ",
                                         average_quality);
    // report average quality for one input layer

     dfm_save_data();
}


ca_defect_range_s : newtype struct of {
   count : integer;
   min : double;
   max : double;
};
ca_defect_density_s : newtype struct of {
   na : integer;            // dummy
   fall_power : double;
   conversion_factor : double;
};
ca_config_s : newtype struct of {
   layer : polygon_layer;
   mode : critical_area_mode_e;
   modeling : integer     // dummy
   range : ca_defect_range_s;
   defect_density : ca_defect_density_s;
};

cas : ca_layer_configuration_l = {
   {DIFF,CA_OPEN,MODELING,{10,0.022,2.000},{NA,2.5,1e-8}},
   {POLY,CA_OPEN,MODELING,{10,0.022,2.000},{NA,2.5,1e-8}},
   {METAL1,CA_OPEN,MODELING,{10,0.022,2.000},{NA,2.5,1e-8}},
};

get_defect_sizes : function         // function for defect size calculation
(
   min_size : double,
   max_size : double,
   steps     : integer,
   fall_power : double
) returning defect_list : defect_sizes_list
{
/*
     // put actual defect sizes calculation coded by user
*/
    defect_list = {0.011, 0.013, 0.017, 0.035, 0.073, 0.13, 0.24, 0.33,
  0.495, 0.74, 1.0};
}

di : integer = 0;
foreach (config in cas)
{
    // calculate defect size
    one_list = get_defect_sizes(config.range.min, config.range.max,
  config.range.count,
 config.defect_density.fall_power);
   defect_hash[di] = one_list;

     di = di + 1;
}

output_layers: list of write_layer_map_s = {};


di = 0;
layer_n : integer = 1;
foreach (config in cas)
 {
   outs : list of polygon_layer;

  if(config.mode == CA_SHORT)
  {
       // for SHORT mode connect_db is required
       cdb = connect({{{config.layer}, connect_layer}});
       outs = critical_areas( config.layer, defect_hash[di], config.mode,
 cdb);
  }
  else
       outs = critical_areas( config.layer, defect_hash[di], config.mode);

    for(i = 0 to defect_hash[di].size()-1) {
      g_ca_layers[""+i] = outs[i];

        output_layers.push_back( {outs[i], {layer_n}} );
        layer_n = layer_n + 1;
    }

    g_ca_layers_keys = g_ca_layers.keys();

    layer_dr = defect_hash[di];
    fall_power = config.defect_density.fall_power;
    ca_mode = config.mode;

    layer_tagname = config.layer.tagname;

    {@ ""+config.layer; dfm_features(g_ca_layers, get_area);}
}

// write out CA layers
gds_file = gds_library("CA_out.gds");
write_gds
(
   output_library = gds_file,
   layers = output_layers
);
\end{lstlisting}
